\documentclass[11pt]{article}
\usepackage[a4paper,margin=24mm]{geometry}
\usepackage{amsmath,amssymb,amsthm,mathtools,mathrsfs}
\usepackage[hidelinks,hypertexnames=false]{hyperref}
\setlength{\emergencystretch}{2em}\allowdisplaybreaks
\newtheorem{theorem}{Theorem}\newtheorem{lemma}[theorem]{Lemma}
\newcommand{\C}{\mathbb C}\newcommand{\R}{\mathbb R}
\newcommand{\norm}[1]{\left\lVert#1\right\rVert}
\title{Subexponential heat cost and the original arithmetic currents}
\author{Split-Zero research programme}\date{19 September 2026}
\begin{document}\maketitle
We receive Sections 1--4 and 8 of the supplied continuation
\cite{web}, prove their heat and current estimates in the original
coordinates, and retain the distinction between the exact currents and
the two specified approximations of them. The full root polynomial,
physical measure, fixed period, original boundary vectors and both
integration regions are retained. Equations HG1--20 also carry the new
growth estimate into the preceding NP heat-depth and leakage receivers.

\section{The original finite source and the exact growth constant}
Use $D=-x\partial_x$, $\mathcal H=L^2(\R_{>0},dx)$ and the full divisor
$h(w)=\prod_{\rho\in Z}(w-\rho)^{m_\rho}$ of $g(w)=2\xi(w)$.
Every $m_\rho$ is the full multiplicity. Set $d=\deg h$,
$F_h(y)=(g/h)(1/2+iy)$ and $w_h(y)=|F_h(y)|^2/(2\pi)$.
HM1--7 \cite{HM} constructs $a_h$ with Mellin transform $g/h$ and
the exact divided-jet section $R_hu=r_h(u)(D)a_h$, where
$j_hr_h(u)=j_h(h/g)u$. Thus
\[
 \mathcal A_NP=P(D)a_h,\quad
 \|\mathcal A_NP\|^2=\int_\R|P(1/2+iy)|^2w_h(y)\,dy,
 \quad u=\mathsf T_hj_hP,\quad \mathsf T_h=j_h(g/h)\cdot .\tag{HG1}
\]
The unit $\mathsf T_h$ includes every original divided derivative.
For $N\ge d$, the full coefficient Gram and attained value form are
\[
\begin{gathered}
 \mathbb G_N=\begin{pmatrix}A_0&C_N\\C_N^*&S_N\end{pmatrix},\quad
 A_0=R_h^*R_h,\quad H_N^{(1)}=A_0-C_NS_N^{-1}C_N^*,
 \\
 c_N^2=\lambda_{\max}((H_N^{(1)})^{-1/2}A_0(H_N^{(1)})^{-1/2}).
\end{gathered}
 \tag{HG2}
\]
The final identity follows directly from HM10: insert
$A_0=H_N^{(1)}+C_NS_N^{-1}C_N^*$ and take the largest eigenvalue.
The attained form is the minimum of (HG1) at fixed $u$, over the complete
relation ideal $h\mathcal P_{N-d}$. This is the quadratic minimum of
\cite[Appendix C.4.1]{Boyd}, proved in HM by completing the square.

\begin{theorem}[Subexponential growth with all density zeros retained]
On this fixed-divisor domain, $\log c_N=o_h(N)$ as $N\to\infty$.
For each fixed admissible $R$, the explicit estimate (HG5) below holds.
\end{theorem}
\begin{proof}
First consider the stipulated quartet. Its midpoint is not a zero of
$g$: the convergent alternating series for $\eta(1/2)$ is a sum of
positive consecutive pairs, hence positive, and
$\zeta(1/2)=\eta(1/2)/(1-\sqrt2)\ne0$. Also $h(1/2)\ne0$.
Define $f_0=|F_h(0)|$, $M_0=\max_{|z|=1}|F_h(z)|$,
$b_0=\min(1/2,f_0/(4M_0))$ and $w_*=f_0^2/(8\pi)$.
Cauchy's coefficient estimate gives
$|F_h(y)-F_h(0)|\le M_0b_0/(1-b_0)\le f_0/2$ on $[-b_0,b_0]$;
therefore $w_h\ge w_*$ there.

For Legendre polynomials the recurrence recorded in
\cite[Table 18.9.1]{DLMF18} proves
$|P_n(z)|\le(1+2r)^n$ on $|z|\le r$. Indeed the induction uses
$2r/(1+2r)+1/(1+2r)^2\le1$, and the initial two polynomials obey
the bound. Their squared norm on $[-1,1]$ is $2/(2n+1)$.
Expansion in the orthonormal basis on $[-b_0,b_0]$ and
Cauchy--Schwarz consequently proves, with
$R_h^{\rm root}=\max_{h(\rho)=0}|(\rho-1/2)/i|$ and
$B_h=1+2(R_h^{\rm root}+1)/b_0$,
\[
 \sup_{|z|\le R_h^{\rm root}+1}|P(1/2+iz)|
 \le\frac{(N+1)B_h^N}{\sqrt{2b_0w_*}}\|\mathcal A_NP\|.
 \tag{HG3}
\]
Here $\sum_{n=0}^N(2n+1)=(N+1)^2$. Cauchy's formula on each
unit circle about an original jet point bounds every divided derivative
by the same supremum. The coordinate change $w=1/2+iz$ contributes
only powers of $i$ of modulus one, which are retained in $j_h$.
Taking the minimum at fixed $u$ yields
\[
 H_N^{(1)}\succeq
 \frac{2b_0w_*}{d\|\mathsf T_h\|^2(N+1)^2B_h^{2N}}I,
 \qquad c_N\le C_h(N+1)B_h^N,
 \quad C_h=\sqrt{\frac{d\|\mathsf T_h\|^2\|A_0\|}{2b_0w_*}}.
 \tag{HG4}
\]

For the sharper rate take $R>0$ with neither endpoint a real zero
of the entire function $F_h$. Its real zeros in this interval are
finite. Let $p_R(y)=\prod_{F_h(t)=0,\ |t|\le R}(y-t)^{n_t}$,
with full analytic multiplicities, $d_R=\deg p_R$, and
$\omega_R=\min_{[-R,R]}|F_h/p_R|^2/(2\pi)>0$.
Choose fixed small circles about the jet points on which $F_h$ has no
zero, possible because $g/h$ is a unit at every full root of $h$.
Let their radii be $r_\rho$, let $R_0$ bound their distance from zero
in the $y$ coordinate, and put
$m_R=\min_{\text{circles}}|p_R|>0$ and
$D_h^2=\sum_{\rho}\sum_{j<m_\rho}r_\rho^{-2j}$.
The full density gives
$\|p_R P(1/2+i\cdot)\|_{L^2[-R,R]}
\le\omega_R^{-1/2}\|\mathcal A_NP\|$.
Apply the preceding Legendre proof to this degree-$(N+d_R)$
polynomial on $[-R,R]$, and divide by $p_R$ only on those circles.
Cauchy's formula and the original $\mathsf T_h$ give
\[
 \begin{gathered}
 c_N\le C_{h,R}(N+d_R+1)(1+2R_0/R)^{N+d_R},\\
 C_{h,R}=\frac{\sqrt{\|A_0\|}\,\|\mathsf T_h\|D_h}
                    {m_R\sqrt{2R\omega_R}}<\infty .
 \end{gathered}\tag{HG5}
\]
For each fixed $R$, divide its logarithm by $N$ and let $N$ tend
to infinity. The upper limit is at most $\log(1+2R_0/R)$.
There are arbitrarily large admissible $R$. Letting them tend to
infinity proves the upper limit zero. Since $H_N^{(1)}\preceq A_0$,
(HG2) gives $c_N\ge1$ and the lower limit is at least zero.
All constants depending on $R$ were fixed before the $N$ limit.
\end{proof}
For any family of one-factor degrees at most $2q+2$ the same bound
also proves $\log c_*=o_h(q)$, where $c_*\ge1$ is their maximum.
Indeed (HG5) bounds all of them by its value at $2q+2$ before
taking the same successive limits. The preceding HM divergence
$c_N\to\infty$ is compatible with this zero exponential rate.

\section{Exact currents before approximating their metric}
At the original cutoff retain the positive $G=G_N$, actual observation
$\Lambda:E\to\operatorname{im}\Lambda$, and
\[
 Q=(\Lambda G^{-1}\Lambda^*)^{-1},\quad
 \Omega=\Lambda^*Q\Lambda,\quad L=G-\Omega=GP_K.\tag{HG6}
\]
Keep $a_1=b_N,a_2=b_{N+1},v=v_\lambda$, their original scale
$\omega_N>0$, $F_i=a_i^*Gv$, $K_i=a_i^*Lv$,
$B_i=a_i^*\Omega v$, $E=v^*Gv$, $d_i=a_i^*Ga_i$ and
$\kappa_N=\sqrt{d_1d_2}/\omega_N$.
ORE/FW \cite{FW} gives
$\mathfrak c_N=\overline{F_1}F_2/(\omega_NE)$, with
$U=K_1/F_1,V=K_2/F_2$. Substitution into the original products
$\overline UV,(1-\overline U)(1-V)$ and their complementary sum gives
\[
 \begin{aligned}
 \Phi_{K,N}&=2\Re(\overline{K_1}K_2)/\omega_N,\\
 \Phi_{B,N}&=2\Re(\overline{B_1}B_2)/\omega_N,\\
 \Phi_{\times,N}&=2\Re(\overline{K_1}B_2+\overline{B_1}K_2)/\omega_N.
 \end{aligned}\tag{HG7}
\]
This is an equality on the original nonzero-denominator domain and
a polynomial continuation across vanishing approximate $F_i$.
Both complementary cross products remain present.

Use the same vectors and $\omega_N$ for $G'$ satisfying
$(1-\eta)G\preceq G'\preceq(1+\eta)G$, $0<\eta\le1/4$.
Taking the minimum on each unchanged observation fibre proves the
same inequalities for $Q$, hence for $\Omega$. Conjugation by
$G^{-1/2}$ gives
$|K_i'-K_i|\le2\eta\sqrt{d_iE}$ and
$|B_i'-B_i|\le\eta\sqrt{d_iE}$, exactly as HR3 \cite{HR}.
Also $|K_i|,|B_i|\le\sqrt{d_iE}$ because $0\preceq L,\Omega\preceq G$.
Expanding all linear and quadratic terms in (HG7) proves
\[
 \begin{aligned}
 |\widetilde\Phi_K-\Phi_K|&\le2\kappa_NE(4\eta+4\eta^2),\\
 |\widetilde\Phi_B-\Phi_B|&\le2\kappa_NE(2\eta+\eta^2),\\
 |\widetilde\Phi_\times-\Phi_\times|&\le2\kappa_NE(6\eta+4\eta^2).
 \end{aligned}\tag{HG8}
\]
Since $|E'-E|\le\eta E$, $E'\ge(1-\eta)E$, and the three
original absolute currents divided by $2E$ are at most
$\kappa_N,\kappa_N,2\kappa_N$, respectively,
\[
 \left|\frac{\widetilde\Phi_{X,N}}{2E'}-
              \frac{\Phi_{X,N}}{2E}\right|\le12\kappa_N\eta,
 \qquad X=K,B,\times.\tag{HG9}
\]
For example the cross bound is
$\kappa_N(8\eta+4\eta^2)/(1-\eta)\le12\kappa_N\eta$;
the other two are smaller. On the nonzero approximate denominator
domain, HR's different approximation keeps $\mathfrak c_N$ fixed.
With $\mathfrak c_N'=\overline{F_1'}F_2'/(\omega_NE')$ its exact
relation to (HG7) is
\[
 \widetilde\Phi_{X,N}-\Phi^{\rm HR}_{X,N}
       =2E'\Re[(\mathfrak c_N'-\mathfrak c_N)P_X'] .\tag{HG10}
\]
Thus the two approximations are connected by a computed correction,
and are not silently identified. In particular their complete sums
retain this correction as well.

\section{Whole-window heat depth and complete receiving error}
On the simple quartet $q=(k+1)^2$ the existing FW48--50 estimates
hold uniformly for $N=q-1+r$, $0\le r\le q+1$:
\[
 \begin{gathered}
 -\log p_{1,N}=2q\psi(r/q)+O_h(k+\log q),\quad
 |\log p_{2,N}|=O_h(\log q),\\
 |U_N|\le e^{O_{h,\varpi}(k\log q)},\quad |V_N|\le C_hq/k,\qquad
 \kappa_N^2p_{1,N}p_{2,N}=|\mathfrak c_N|^2 .
 \end{gathered}\tag{HG11}
\]
The original signed correction in FW28 is
$\mathfrak c_N=k\delta+i(-1)^r(k\gamma-\varepsilon_N)$,
$0<\varepsilon_N=O_h(k^3/q^2)$, so
$\log|\mathfrak c_N|=O_h(\log q)$.
The profile is strictly decreasing: FW44 gives
$\psi'(t)=\log(\kappa_t/E(\kappa_t))<0$ for $t>0$,
and $\psi(0)=a_0=\log(4/\pi)$.
For an explicit check, use $r_tK/E=1/(1+t)$ and
$\psi(t)=\log(1+r_t)-(1+t)\log E+(t/2)\log(1-r_t^2)$.
Differentiation using Carlson's formulas \cite{Carlson} cancels
the coefficient of $r_t'$ exactly and leaves the stated derivative.
Here $E\ge1$ and $0<\kappa_t<1$ prove its sign.
Therefore, with $p_*^{\rm win}=\min_{N,i}p_{i,N}$,
\[
 -\log p_*^{\rm win}=2a_0q+O_h(k+\log q),\qquad
 \log\max_N\kappa_N=a_0q+O_h(k+\log q).\tag{HG12}
\]

HM18--19 gives the relative metric comparison with
$t_J=2^{-(J+1)/2}c_*$ and
$\eta_J=(1+t_J)^{2k}-1$, provided $t_J<1$.
For a specified $0<\eta_*\le1/4$, a finite successful depth is
\[
 J=\max\{0,\lceil2\log_2(\bar c_k/d_*)\rceil-1\},\quad
 d_*=(1+\eta_*)^{1/(2k)}-1,\quad
 \bar c_k=C_h(2q+3)B_h^{2q+2}.\tag{HG13}
\]
The least common strict HR denominator guard is
$\eta_J<\sqrt{p_*^{\rm win}}$. Inverting this increasing scalar
function and retaining its strict integer rounding yields
\[
 \begin{gathered}
 J_{\rm guard}+1=2\log_2c_*+2\log_2(2k)
                      -\log_2p_*^{\rm win}+O(1),\\
 \displaystyle\lim_{k\to\infty}\frac{J_{\rm guard}}q
                         =\frac{2\log(4/\pi)}{\log2}.
 \end{gathered}\tag{HG14}
\]
Indeed for $0<x\le1$, $x/2\le\log(1+x)\le x$ and
$b\le e^b-1\le e^{1/2}b$ for $b=\log(1+x)/(2k)$.
These give a bounded uniform error when taking logarithms of
$d_*$. Integer rounding adds at most one. Insert (HG12) and
$\log c_*=o(q)$ to get the limit.

For any fixed $\epsilon_0>0$, set
$J_k=\lceil[2a_0/\log2+\epsilon_0]q\rceil$.
Then $\eta_{J_k}=\exp[-a_0q-(\epsilon_0\log2/2)q+o_h(q)]$.
HR6 gives
$|U'-U|\le\eta(2+|U|)/(\sqrt{p_1}-\eta)$ and the identical
formula for $V$. Substitution of (HG11)--(HG12), and $k\log q=o(q)$,
proves
\[
 \begin{aligned}
 &\max_N(|U_N^{[J_k]}-U_N|+|V_N^{[J_k]}-V_N|)\\
 &\quad+\sum_{N=q-1}^{2q-1}w_N
 \left|\frac{\widetilde\Phi_{X,N}^{[J_k]}}{2E_N^{[J_k]}}
                  -\frac{\Phi_{X,N}}{2E_N}\right|
 \le e^{-(\epsilon_0\log2/2)q+o_{h,\varpi}(q)} .
 \end{aligned}\tag{HG15}
\]
The weights are the original $1,2,\ldots,2,1$, whose sum is $2q$.
Their contribution has logarithm $O(\log q)$ and is absorbed in the
displayed remainder. This controls approximation of each full signed
current. It does not supply its unevaluated exact sign.
For a determinant return of total absolute weighted rank $B_{\det}$,
the same minimum comparison instead gives
$|\mathcal L'-\mathcal L|\le B_{\det}[-\log(1-\eta)]$.
The original polynomial-rank returns therefore require only
$J=2\log_2c_*+O(\log q+\log^+(1/\varepsilon))=o(q)$ at fixed tolerance.

\section{The arithmetic action and its original boundary}
Let $r=\operatorname{rem}_hP$ and
$\ell_h(P)=[w^{d-1}]r$. In the literal polynomial coordinate the
heat source is
$\mathcal A_N^{[J]}P=(I-T_s)P(D)a_h+T_sr(D)a_h$,
$T_sf(x)=f(sx)$, $s=2^{J+1}$.
Euclidean division gives
$wr-\operatorname{rem}_h(wP)=h\ell_h(P)$.
Since $D$ commutes with $T_s$ and $h(D)a_h=f_0=\Theta\phi_*$,
\[
 D\mathcal A_N^{[J]}P-\mathcal A_{N+1}^{[J]}(wP)
             =T_sf_0\ell_h(P),\qquad T_sf_0=\Theta(T_s\phi_*).
 \tag{HG16}
\]
The original $\phi_*=(4\pi^2x^4-6\pi x^2)e^{-\pi x^2}$ is even,
Schwartz, vanishes at zero and has zero integral. Dilation preserves
these properties, so the displayed defect belongs to the same theta
source. The fixed Hermite interpolation matrix sends $j_hP$ to
$\ell_h(P)$. The proof of (HG5), with this fixed row in place of
$\mathsf T_h$, therefore gives
$\log^+\|\ell_h\|_{\mathcal H_N^*}=o_h(N)$.
Here $\mathcal H_N$ is exactly the polynomial form in (HG1).

The original multinomial map is
$P(S)\mapsto P(w_1+\cdots+w_k)$. Differentiate the tensor source
in each physical variable and apply (HG16) in that factor.
Its error is the sum of all $k$ tensor terms, each with one factor
$T_sf_0\ell_h$ and the remaining original heat source factors.
HM's isometric original-source norm and $\|T_s\|=s^{-1/2}$ give
\[
 \|\mathfrak D_{k,N}^{[J]}\|
 \le k s^{-1/2}\|f_0\|\|\ell_h\|_{\mathcal H_N^*}
                  (1+t_J)^{k-1}.\tag{HG17}
\]
Restricting through the multinomial map has norm one when its domain
uses its pullback, the full original convolution form. Thus this
bound keeps that norm, rather than an unrelated coefficient norm.
At $J=\lceil\epsilon_0q\rceil$ and $N\le2q+2$ it is
$\exp[-(\epsilon_0\log2/2)q+o_h(q)]$.

Let $\mathsf T_N:E_\chi\to\mathcal P_N$ be the original attained
section for the full root polynomial, and $A=M_S$ on $E_\chi$.
Adding and subtracting the same polynomial lift proves exactly
\[
 \begin{aligned}
 D_{\rm tot}\mathcal A_{k,N}^{[J]}\mathsf T_N
 -\mathcal A_{k,N+1}^{[J]}\mathsf T_{N+1}A
 ={}&\mathfrak D_{k,N}^{[J]}\mathsf T_N\\
 &+\mathcal A_{k,N+1}^{[J]}
       (M_S\mathsf T_N-\mathsf T_{N+1}A).
 \end{aligned}\tag{HG18}
\]
The parenthesized map has zero original remainder and takes values
in $\chi\mathcal P_{N+1-q}$. Its entire image, and the heat term,
are retained in the complete arithmetic action.

\section{Back-propagation into the established native receivers}
The full original observability theorem O1--12 \cite{native} proves
$\Lambda v_\lambda\ne0$ on its explicit period locus. Consequently
$e_{B,N}=(\Lambda v)^*Q_N\Lambda v=E_N(1-\theta_N)>0$ there.
FW39 and NP19 \cite{NP} therefore apply without an undetected-class
alternative. For heat constants $\alpha=(1-t_J)^{2k}$,
$\beta=(1+t_J)^{2k}$, minimum comparison gives
$\alpha/\beta\le(1-\theta_N')/(1-\theta_N)\le\beta/\alpha$.
The error in the original four-endpoint angle return is at most
\[
 4\log(\beta/\alpha)=8k\log\frac{1+t_J}{1-t_J}.\tag{HG19}
\]
There is no division by the possibly small observed fraction.
Likewise the NP leakage Gram uses the same fixed action blocks in
the original $G_N$; its proved quantitative estimate stays
\[
 \sum_{j=0}^{r-1}(T^j)^*B^*Q_NBT^j
 \succeq(C_{R,N}C_{W,N}C_{T,N})^{-2}H_{K,N}.\tag{HG20}
\]
The new heat-growth theorem supplies finite approximation control for
the metrics entering these original objects. It does not change the
proper-source quotient into the different source-endpoint quotient,
or discard the target minimum in NP30--31. The remaining native
allocation and current values must still be evaluated in those maps.

\section{A finite common guard with every fixed primary multiplicity}
The sharp coefficient (HG14) uses the simple-packet profile. For the
full packet $q=[1+k(m_0-1)](k+1)^2$, the incoming finite guard can
also be proved without an observation inverse. Retain GRA3--6
\cite{GRA}: put $p_h=8m_0-9/2$, $B_h^{\rm tail}=42+8m_0$,
$M_h^{\rm tail}=21+4m_0$, $\alpha_* =\pi/2$ and
\[
\begin{gathered}
\vartheta_h=\int_{-1}^1w_h(y)\,dy,\quad
M_h(t)=\int_\R e^{ty}w_h(y)\,dy,\quad
c_\Gamma=\sqrt2e^{-7/6},\quad C_\Gamma=\sqrt{2\sqrt5}\,e^{2/3},\\
\underline a_k=\frac{c_h\vartheta_h^{k-3}e^{-\alpha_*(k-3)}
 (k-2)^{-B_h^{\rm tail}}}{C_\Gamma2^{B_h^{\rm tail}/2}},\quad
\overline a_k=\frac{C_h^{\rm tilt}}{c_\Gamma}
 k^{p_h+1}M_h(\alpha_*)^{k-1},\\
D_n=[1+4(n+M_h^{\rm tail})^2]^{M_h^{\rm tail}},\qquad
\frac{\underline a_k}{D_n}\|P\|_\sigma^2
\le\|P\|_{w_h^{*k}}^2\le\overline a_k\|P\|_\sigma^2
\quad(\deg P\le n).
\end{gathered}\tag{HG21}
\]
Here $c_h,C_h^{\rm tilt}$ are precisely the full-source TW7 and BT12
constants retained in GRA4. The monic Gamma norm in this exact measure
is $\gamma_n=\sqrt{2\pi}(2n)!/4^n$.
For clarity GRA's polynomial comparison has no degree restriction:
its recurrence $\|yP\|_\sigma\le2(n+1)\|P\|_\sigma$ gives
$\int|P|^2(1+y^2)^{M_h^{\rm tail}}d\sigma\le D_n\|P\|_\sigma^2$.
Cauchy--Schwarz and the complete density lower bound
$w_h^{*k}\ge\underline a_k(1+y^2)^{-M_h^{\rm tail}}\sigma$
give the first inequality in (HG21); the full upper density gives
the second. These retain all normalization masses.

Set $n_*=2q+2$, $R_k=k\sqrt{\delta^2+\gamma^2}$ and $\tau=\pi/4$.
For the complete monic root polynomial $Q$ in $S=k/2+iy$ all roots
have modulus at most $R_k$, with every multiplicity retained.
The literal monic trial $y^rQ(y)$ for the rank-$r$ relation minimum
$\nu_r$ obeys
\[
\nu_r\le\int(|y|+R_k)^{2(q+r)}w_h^{*k}(y)\,dy
 \le2e^{\tau R_k}M_h(\tau)^k(2q+2r)!\tau^{-2q-2r}.
\]
The second inequality uses $t^{2n}\le(2n)!\tau^{-2n}e^{\tau t}$
and evenness of the full density. The original source minimum
satisfies $\omega_n\ge\underline a_k\gamma_n/D_{n_*}$ for
$n\le n_*$. Therefore
\[
\begin{gathered}
\Delta_r=\omega_{q+r}/\nu_r\ge\underline\Delta_k,
\quad\underline\Delta_k=
\min\left\{\frac12,
\frac{\underline a_k\sqrt{2\pi}}
{2D_{n_*}e^{\tau R_k}M_h(\tau)^k}
\left(\frac{\tau^2}{4}\right)^{n_*}\right\},\\
\overline b_k=\max\{1,
(\overline a_kD_{n_*}/\underline a_k)(n_*+1)^2\},\quad
\kappa_N\le\overline\kappa_k=
\sqrt{\overline b_k/\underline\Delta_k},\qquad
p_{i,N}\ge\min\{1/2,k^2\delta^2/\overline\kappa_k^2\}=:p_*.
\end{gathered}\tag{HG22}
\]
Indeed (HG21) bounds $\omega_{n+1}/\omega_n$ by
$\overline b_k$. The original FW boundary-radius formula is
$\kappa_{q-1}^2=(\omega_q/\omega_{q-1})(\Delta_0^{-1}-1)$ and,
for $r=N+1-q\ge1$,
$\kappa_N^2=(\omega_{N+1}/\omega_N)(1-\Delta_{r-1})(\Delta_r^{-1}-1)$.
Every extra factor is at most its positive upper bound in (HG22).
Finally $\kappa_N^2p_{1,N}p_{2,N}=|\mathfrak c_N|^2\ge k^2\delta^2$
and $p_{i,N}\le1$ prove both denominator bounds.
GRA6 gives $\log(\overline a_k/\underline a_k)=O_h(k+\log q)$;
each remaining logarithm in (HG22) is $O_h(q)$, hence
$\log\overline\kappa_k+\log(1/p_*)=O_h(q)$.

For any prescribed $\varepsilon>0$, choose
\[
\eta_*\le\min\left\{\frac14,\frac{\varepsilon}{32q\overline\kappa_k},
\frac{\sqrt{p_*}}2,\frac{\varepsilon p_*}{6},
\frac{\varepsilon}{2(B_{\det}+\varepsilon)}\right\} .\tag{HG23}
\]
Then (HG13) gives weighted current error at most $\varepsilon$,
each response error at most $\varepsilon$ by HR6--7, and determinant
return error at most $\varepsilon$, using
$-\log(1-\eta_*)\le2\eta_*$.
For the original polynomial-rank returns the resulting dyadic depth
is $O_h(q+\log^+(1/\varepsilon))$. All finite guards are explicit;
this is not a bound on the computational complexity of the integrals.

\begin{thebibliography}{10}
\bibitem{web} Supplied web continuation, \emph{GitHub integration and further
calculations}, 19 September 2026, equations (1)--(31), (63)--(64).
Full mathematical response preserved in \texttt{RECEIVED\_MATHEMATICS.md}.
\bibitem{HM} \emph{The circle heat source in the original arithmetic quotient
metrics}, HM1--30; full provider included, public commit
\texttt{607728ee3c35be98720d548794727670b1946c91}.
\bibitem{HR} \emph{Heat approximation errors for the original complex
observation responses}, HR1--11, same public edition; full provider included.
\bibitem{FW} \emph{Full-window source products and the retained class},
FW1--52, public commit \texttt{3c44de632b2c5ead13b23619d646803d2bbffc19};
full provider included.
\bibitem{GRA} \emph{Growing-degree original Gamma comparison}, GRA1--6,
complete proof in the supplied cumulative signed-return source;
original constants also in \texttt{EC\_COMPLETE.md}, EC14--15.
\bibitem{NP} \emph{Native precision, marked detection and quantitative kernel
leakage}, NP1--31, 19 September 2026; full preceding proof included.
\bibitem{native} \emph{Full original observability}, O1--12, and
\emph{Original conductor coprimality}, NC1--27, public commit
\texttt{12dd3645229bfebbf1cb8f5968ecffd4dfff1e9f}.
\bibitem{DLMF18} T. H. Koornwinder, R. Wong, R. Koekoek, R. F. Swarttouw
and W. P. Reinhardt, \emph{Orthogonal Polynomials}, NIST DLMF Chapter 18,
\href{https://dlmf.nist.gov/18.9.T1}{Table 18.9.1} and
\href{https://dlmf.nist.gov/18.3}{Section 18.3}.
\bibitem{Carlson} Bille C. Carlson, \emph{Elliptic Integrals}, NIST DLMF
Chapter 19, \url{https://dlmf.nist.gov/19.4}.
\bibitem{Boyd} Stephen Boyd and Lieven Vandenberghe,
\emph{Convex Optimization}, Cambridge University Press, 2004, Appendix C.4.1,
\url{https://web.stanford.edu/~boyd/cvxbook/bv_cvxbook.pdf}.
\end{thebibliography}
\end{document}
